\section{Conclusion}
\label{sec:conclusion}

We have presented two parameter-efficient architectural additions to hybrid state-space and attention language models. \emph{Hormone routing} is a token-conditional router that injects a gated combination of frozen direction vectors into the residual stream after each attention block. \emph{Intracellular attention} is a slot-attention mechanism placed inside the SSM scan that lets the otherwise-independent state slots interact via content-dependent attention before the linear readout. Both add fewer than $1\%$ of trainable parameters and are gated such that the unmodified hybrid backbone is recovered exactly at initialisation.

Our central empirical finding is a clean decoupling between two quantities that interpretability work typically assumes are correlated. The hormone-routing module develops semantically-coherent per-genre routing specialisation across every variant we examined --- including controls that replace the extracted directions with random unit vectors and that freeze the output gate at $1.0$ --- with maximum pairwise $L_1$ distance between genre routing distributions in the range $0.6$--$1.6$ (uniform $= 0$, disjoint $= 2$). The variant with the \emph{highest} specialisation strength is one that the model cannot modulate (\textsc{HR-fixedgate}, gate frozen at $1.0$). Yet at $125$M parameters and $2$B tokens, none of this organisation produces a measurable contribution to the loss objective: all properly-initialised variants converge to identical validation perplexity. \textbf{Specialisation strength does not correlate with loss.} The forced-engagement variant \TBD{(\textsc{HR-forced} narrative)} provides partial discrimination between two readings of this decoupling: the architecture-and-data system produces representational structure independent of whether that structure is loss-relevant, and at this scale that decoupling is essentially complete.

The mechanism is bounded --- random vectors and frozen gates yield base-level loss --- but sensitive to the training schedule: introducing hormone routing during base pretraining (rather than as a continuation from a stabilised checkpoint) actively degrades performance, even though the router still develops a discriminative routing distribution.

\paragraph{Future directions.} A fused CUDA kernel for intracellular attention would close the gap to the standard Mamba scan and remove the practical barrier to applying the mechanism at scale; we have established the architecture and a preliminary engagement signal, but the full empirical case requires the kernel. The hormone-extraction protocol, currently performed once at the $1$B-token milestone, could be made continuous --- updating $V$ on a slow schedule throughout training --- which would let the bank track the model's evolving representations. Most importantly, the routing-organisation-without-loss-relevance finding should be tested at larger scale: whether the bounded-engagement regime we observe at $125$M parameters reflects a fundamental property of the mechanism or a capacity-bottlenecked phenomenon depends on how the residual stream's dimensionality interacts with the perturbation magnitude. The forced-engagement variant we report begins this investigation but does not settle it.

A broader research programme suggested by this work is the systematic study of \emph{introspective architectural channels}: components that develop structured internal organisation along axes the architecture itself defines, observable to the model and to interpretation, but decoupled --- at least at small scale --- from the optimisation objective. The hormone-routing module is one example. Whether such channels become loss-relevant at scale, and what conditions cause them to do so, remains an open question we have only begun to examine.
